\scanpage{109}

Donc on trouve

\[w_*(f_{V_{n+1}}(\beta_0,\ldots,\beta_{n-1}))=\bigl[1+\underset{\substack{\rotatebox{-90}{$\in$}\\ H^1(-,\deleted{\unsure{\mu_2|\mu}}\mathbb Z/2)}}{\underline{\partial(\widetilde\rho_{0,n-1}\widetilde\rho_{1,n-1})}}\bigr]\,w_*(f_{V_n}(\beta_{\unsure{1|{}}},\ldots,\beta_{n-1}))\]

On trouve ainsi de proche en proche, pourvu que les calculs aient un sens i.e. que \deleted{tous} les $\rho_{i,n-1}$ soient inversibles $(0\leq i\leq n-1)$

\[\begin{aligned}w_*(f_{V_{n+1}}(\beta_0,\ldots,\beta_{n-1}))={}&[1+\partial(\widetilde\rho_{0,n-1}\widetilde\rho_{1,n-1})][1+\partial(\widetilde\rho_{1,n-1}\widetilde\rho_{2,n-1})]\\ &\ldots[1+\partial(\widetilde\rho_{\deleted{\unsure{n-2,n-1}}\inserted{n-2,n-1}}\widetilde\rho_{\deleted{\unsure{n-2,n-1}}\inserted{n-1,n-1}})]\,\unsure{\bar w_*|w_*}(f_{V_2}(\beta_{n-1}))\end{aligned}\]

De $=$, on trouve,

\[\begin{aligned}w_*(f_{V_2}(\beta_{n-1}))&=\deleted{\bigl(1+\partial[\underbrace{\widetilde\rho_{n-1,n-1}}_{4-\beta_{n-1}}\underbrace{\widetilde\rho_{n,n-1}}_{\Vert\,1}]\bigr)\underbrace{\bigl(1+\partial(\underbrace{\widetilde\rho_{n,n-1}}_{1}\underbrace{\widetilde\rho_{n+1,n}}_{1})\bigr)}_{\unsure{1}}}\\ &\bigl(1+\partial(\beta(4-\beta))\bigr)=[1+\partial(2+\alpha)\deleted{\unsure{+}}(2-\alpha))]\end{aligned}\]

En effet, faisant le calcul pour $f_V(\beta)$ sur $a_0,a_1$ on trouve, si on veut

\[\left\{\begin{aligned}f(a_0)&=1\qquad\partial f(a_0)=0\\a_0^\perp&=\{x\unsure{,|{}}a_0+ya_1\mid 2x-\beta y=0\quad\text{i.e. }\deleted{y=}x=\tfrac\beta2y\}\end{aligned}\right.\]

engendré par $\beta a_0+2a_1=u$

\[\begin{aligned}f(u)&=\deleted{\unsure{\beta^2+4\beta-|\beta^2+4\beta}}\ \beta^2\underbrace{f(a_0)}_1+4\underbrace{f(a_1)}_\beta+2\beta\underbrace{\varphi(a_0,a_1)}_{-\beta}\\ &=4\beta-\beta^2=\beta(4-\beta)\end{aligned}\]

d'où \deleted{$w_*($} la formule annoncée.

% archive sheet number: 108
